\documentclass[10pt,a4paper]{amsart}
\usepackage{amsmath,amssymb,amsthm,graphicx}
% For commutative diagrams we use Xy-pic
\usepackage[all]{xy}

% We use 2cell for 2-commutative diagrams.
\xyoption{2cell}
% Theorem environments.
%
\theoremstyle{plain}
\newtheorem{theorem}[subsection]{Theorem}
\newtheorem{proposition}[subsection]{Proposition}
\newtheorem{lemma}[subsection]{Lemma}

\theoremstyle{definition}
\newtheorem{definition}[subsection]{Definition}
\newtheorem{example}[subsection]{Example}
\newtheorem{exercise}[subsection]{Exercise}
\newtheorem{situation}[subsection]{Situation}

\theoremstyle{remark}
\newtheorem{remark}[subsection]{Remark}
\newtheorem{remarks}[subsection]{Remarks}

\numberwithin{equation}{subsection}

% Macros
%
\def\lim{\mathop{\mathrm{lim}}\nolimits}
\def\colim{\mathop{\mathrm{colim}}\nolimits}
\def\Spec{\mathop{\mathrm{Spec}}}
\def\Hom{\mathop{\mathrm{Hom}}\nolimits}
\def\Ext{\mathop{\mathrm{Ext}}\nolimits}
\def\SheafHom{\mathop{\mathcal{H}\!\mathit{om}}\nolimits}
\def\SheafExt{\mathop{\mathcal{E}\!\mathit{xt}}\nolimits}
\def\Sch{\mathit{Sch}}
\def\Mor{\mathop{\mathrm{Mor}}\nolimits}
\def\Ob{\mathop{\mathrm{Ob}}\nolimits}
\def\Sh{\mathop{\mathit{Sh}}\nolimits}
\def\NL{\mathop{N\!L}\nolimits}
\def\CH{\mathop{\mathrm{CH}}\nolimits}
\def\proetale{{pro\text{-}\acute{e}tale}}
\def\etale{{\acute{e}tale}}
\def\QCoh{\mathit{QCoh}}
\def\Ker{\mathop{\mathrm{Ker}}}
\def\Im{\mathop{\mathrm{Im}}}
\def\Coker{\mathop{\mathrm{Coker}}}
\def\Coim{\mathop{\mathrm{Coim}}}

% Boxtimes
%
\DeclareMathSymbol{\boxtimes}{\mathbin}{AMSa}{"02}

%
% Macros for moduli stacks/spaces
%
\def\QCohstack{\mathcal{QC}\!\mathit{oh}}
\def\Cohstack{\mathcal{C}\!\mathit{oh}}
\def\Spacesstack{\mathcal{S}\!\mathit{paces}}
\def\Quotfunctor{\mathrm{Quot}}
\def\Hilbfunctor{\mathrm{Hilb}}
\def\Curvesstack{\mathcal{C}\!\mathit{urves}}
\def\Polarizedstack{\mathcal{P}\!\mathit{olarized}}
\def\Complexesstack{\mathcal{C}\!\mathit{omplexes}}
% \Pic is the operator that assigns to X its picard group, usage \Pic(X)
% \Picardstack_{X/B} denotes the Picard stack of X over B
% \Picardfunctor_{X/B} denotes the Picard functor of X over B
\def\Pic{\mathop{\mathrm{Pic}}\nolimits}
\def\Picardstack{\mathcal{P}\!\mathit{ic}}
\def\Picardfunctor{\mathrm{Pic}}
\def\Deformationcategory{\mathcal{D}\!\mathit{ef}}
\usepackage{cleveref}
\crefname{originallemma}{original lemma}{original lemmas}
\Crefname{originallemma}{Original lemma}{Original lemmas}
\makeatletter
\newlabel{derived-lemma-third-isomorphism-triangle}{{014A}{source}{}{sourcepointer-1}{}}
\newlabel{derived-lemma-third-isomorphism-triangle@cref}{{[originallemma][0][]014A}{[1][0][]source}{}{}{}}
\makeatother
\crefname{originalreference}{original reference}{original references}
\Crefname{originalreference}{Original reference}{Original references}
\makeatletter
\newlabel{lemma-finite-cd-mod-n-direct-sums}{{0F11}{source}{}{sourcepointer-2}{}}
\newlabel{lemma-finite-cd-mod-n-direct-sums@cref}{{[originalreference][0][]0F11}{[1][0][]source}{}{}{}}
\makeatother
\crefname{originalremark}{original remark}{original remarks}
\Crefname{originalremark}{Original remark}{Original remarks}
\makeatletter
\newlabel{more-algebra-remark-P-resolution}{{09PB}{source}{}{sourcepointer-3}{}}
\newlabel{more-algebra-remark-P-resolution@cref}{{[originalremark][0][]09PB}{[1][0][]source}{}{}{}}
\makeatother
\crefname{originalsection}{original section}{original sections}
\Crefname{originalsection}{Original section}{Original sections}
\makeatletter
\newlabel{sites-cohomology-section-projection-formula}{{0943}{source}{}{sourcepointer-4}{}}
\newlabel{sites-cohomology-section-projection-formula@cref}{{[originalsection][0][]0943}{[1][0][]source}{}{}{}}
\makeatother
\title{Global sections of a derived tensor product with a constant complex pull the constant factor out on a scheme of finite cohomological dimension}
\author{The Stacks Project Authors}
\date{stacks-project:etale-cohomology.tex:local-snapshot-20261011; GNU FDL 1.2 or later}
\begin{document}
\maketitle
\noindent\textit{Editorial note.} All delivered mathematical text is the author's own, copied from the pinned Stacks Project chapter file etale-cohomology.tex; no statement or proof step was rewritten, reordered or removed.\par\medskip
\noindent\textit{Editorial note.} The counted proof is the author's complete argument: the proof is written out in full: the canonical map, the property T and the three conditions of the resolution criterion, with the reason for each condition and the base case.\par\medskip
\noindent\textit{Editorial note.} The label lemma-pull-out-constant-mod-n names the counted statement, and the proof environment immediately following it in the source is delivered verbatim as its proof.\par\medskip
\noindent\textit{Editorial note.} The counted statement is carried by the author's own lemma environment in the chapter source.\par\medskip
\noindent\textit{Editorial note.} The statement and proof are copied from the pinned chapter file; the four cited inputs are kept as external pointers with their canonical Stacks tags.\par\medskip
\noindent\textit{Allowed standard prerequisites (editorial).}\begin{itemize}
\item The canonical projection map. Cohomology on Sites, the section on the projection formula providing the canonical map R Gamma(X, E) tensor\textasciicircum{}L K -> R Gamma(X, E tensor\textasciicircum{}L underline K). Supplies the map whose being an isomorphism is proved.
\item Commutation with direct sums for finite cohomological dimension. Same chapter, the lemma stating that the cohomology of a scheme of finite cohomological dimension commutes with direct sums in the relevant derived category. Supplies property (1): both sides commute with direct sums.
\item The resolution criterion for complexes over a ring. More on Algebra, the remark listing the three conditions sufficient for a property of complexes to hold for all K, together with Derived Categories, the lemma comparing exact functors on triangles. Supplies conditions (2) and (3) of the criterion and its conclusion.
\end{itemize}\medskip
\ifdefined\hypertarget\hypertarget{sourcepointer-1}{}\fi
\noindent\textit{Source reference.} Original lemma 014A: sources/stacks-chow/tags.txt, source lines 1465--1465. Prerequisite pointer invoked at proof line(s) 20187; it names a standard fact of a sibling chapter that the counted argument uses, and every step of the counted proof built on it is retained verbatim in the delivered body.\par\medskip
\ifdefined\hypertarget\hypertarget{sourcepointer-2}{}\fi
\noindent\textit{Source reference.} Original reference 0F11: sources/stacks-etale-cohomology/pinned-original-source.tex, source lines 20074--20074. Prerequisite pointer invoked at proof line(s) 20184; it names a standard fact of the same chapter that the counted argument uses, and every step of the counted proof built on it is retained verbatim in the delivered body.\par\medskip
\ifdefined\hypertarget\hypertarget{sourcepointer-3}{}\fi
\noindent\textit{Source reference.} Original remark 09PB: sources/stacks-chow/tags.txt, source lines 12244--12244. Prerequisite pointer invoked at proof line(s) 20180; it names a standard fact of a sibling chapter that the counted argument uses, and every step of the counted proof built on it is retained verbatim in the delivered body.\par\medskip
\ifdefined\hypertarget\hypertarget{sourcepointer-4}{}\fi
\noindent\textit{Source reference.} Original section 0943: sources/stacks-chow/tags.txt, source lines 11536--11536. Prerequisite pointer invoked at proof line(s) 20176; it names a standard fact of a sibling chapter that the counted argument uses, and every step of the counted proof built on it is retained verbatim in the delivered body.\par\medskip
\begin{lemma}
\label{lemma-pull-out-constant-mod-n}
Let $X$ be a quasi-compact and quasi-separated scheme
such that $\text{cd}(X) < \infty$. Let $\Lambda$ be a torsion ring.
Let $E \in D(X_\etale, \Lambda)$ and $K \in D(\Lambda)$. Then
$$
R\Gamma(X, E \otimes_\Lambda^\mathbf{L} \underline{K}) =
R\Gamma(X, E) \otimes_\Lambda^\mathbf{L} K
$$
\end{lemma}

\begin{proof}
There is a canonical map from right to left by
Cohomology on Sites, Section \ref{sites-cohomology-section-projection-formula}.
Let $T(K)$ be the property that the statement of the
lemma holds for $K \in D(\Lambda)$.
We will check conditions (1), (2), and (3) of
More on Algebra, Remark \ref{more-algebra-remark-P-resolution}
hold for $T$ to conclude.
Property (1) holds because both sides of the equality
commute with direct sums, see
Lemma \ref{lemma-finite-cd-mod-n-direct-sums}.
Property (2) holds because we are comparing exact functors
between triangulated categories and we can use
Derived Categories, Lemma \ref{derived-lemma-third-isomorphism-triangle}.
Property (3) says the lemma holds when $K = \Lambda[k]$
for any shift $k \in \mathbf{Z}$ and this is obvious.
\end{proof}

\end{document}
